%%%%%%%%%%%%%%%%%%%%%%%%%%%%%%%%%%%%%%%%%%%%%%%%%%%%%%%%%%%%%%%%%%%%%%%%%%%%%%%%
%%% THEORY FOR 'ALKCALC'                                                     %%%
%%%                                                                          %%%
%%% AUTHOR OF THIS FILE: SIMON EUCHNER                                       %%%
%%%%%%%%%%%%%%%%%%%%%%%%%%%%%%%%%%%%%%%%%%%%%%%%%%%%%%%%%%%%%%%%%%%%%%%%%%%%%%%%

\documentclass[fontsize=12pt,
               twoside=false,
               BCOR=0mm,
               DIV=14,
               overfullrule=true, % SET 'TRUE' FOR OVERFULL BOX INDICATORS
               %chapterprefix=false, % NOT NEEDED FOR SCRARTCL
               toc=bibliography,
               headinclude=true,
               footinclude=false]{scrartcl}

%%% FONT
\usepackage{libertine}
\usepackage{libertinust1math}
\usepackage{dutchcal}
\KOMAoption{DIV}{last}
\setkomafont{captionlabel}{\usekomafont{descriptionlabel}}
\usepackage{upgreek,textcomp}
\AfterTOCHead[toc]{\sffamily}

%%% CAPTIONS
\setcapindent{0em}
\setcapwidth[c]{\textwidth}
\setkomafont{caption}{\sffamily}

%%% LANGUAGE
\usepackage[UKenglish]{babel}

%%% HEADER/FONT
\usepackage[headsepline=1pt:head,
            footsepline=1pt:foot]{scrlayer-scrpage}

%%% FOOTNOTES
\setfootnoterule[.5pt]{.5\textwidth}
\deffootnote{1em}{1em}{
    \makebox[1em][l]{\thefootnotemark}
}
\setkomafont{footnote}{\sffamily}
\setlength{\skip\footins}{1.5em}
\setlength{\footnotesep}{1em}

%%% SECTION AND SUBSECTION SETTINGS
\renewcommand{\thesection}{\Roman{section}}
\renewcommand{\thesubsection}{\roman{subsection}}

%%% GRAPHICS
\usepackage[pdftex]{graphicx}
\usepackage{tikz,color,xcolor}
\graphicspath{{./}}

%%% TABLES
\usepackage{booktabs}

%%% ALIGN TEXT
\usepackage{enumerate}

%%% DOCUMENTATION
\usepackage{sdoc}

%%% EQUATION TYPESETTING
\usepackage{amsmath,amssymb,amsthm,mathtools}
\usepackage{physics,bm,slashed,cancel,tcolorbox}
\numberwithin{equation}{section}

%%% HYPERLINKS
\usepackage[unicode=true,
            colorlinks=false, % SET TO 'FALSE' FOR COLOURED REF'S, ELSE 'TRUE'
            citebordercolor=blue,
            citecolor=black,
            urlbordercolor=blue,
            urlcolor=black,
            linkbordercolor=orange,
            linkcolor=black,
            pdfborder={0 0 1},
            pdfborderstyle={/S/U/W 1}]{hyperref}

%%% BIBLIOGRAPHY
\usepackage{csquotes,bib}
\addbibresource{./references.bib}

%%% ACRONYMS
\usepackage{acro}
\DeclareAcronym{fem}{
    short=FEM,
    long=finite element method,
    short-plural=s,
    long-plural=s,
}
\DeclareAcronym{fdm}{
    short=FDM,
    long=finite difference method,
    short-plural=s,
    long-plural=s,
}

%%% CUSTOM COMMANDS
\newcommand{\ee}{\mathrm{e}}
\newcommand{\ii}{\mathrm{i}}
\newcommand{\pii}{\uppi}
\newcommand{\ed}{\downarrow}
\newcommand{\eu}{\uparrow}
\newcommand{\aB}{a_\mathrm{B}}
\newcommand{\me}{m_\mathrm{e}}
\newcommand{\EH}{E_\mathrm{H}}
\newcommand{\kBT}{k_\mathrm{B}T}
\newcommand{\alkcalc}{\texttt{AlkCalc} }
\newcommand{\figref}[2]{Fig.~\hyperref[#1]{\ref{#1}(#2)}}
\newcommand{\tus}{\rule{2mm}{.7pt}}

%%% TESTING
\usepackage{lipsum}

%%% TITLE PAGE
\title{Theory and manual\par\vspace*{5mm}for\par\vspace*{5mm}\alkcalc}%
\setcounter{footnote}{1}%
\author{Simon Euchner\footnote{Electronic Mail: %
        $\langle$\texttt{euchner.se@gmail.com}$\rangle$}%
       }
\date{20.08.2025}

%%% TABLES
\def\TabAngularFactors{
    \begin{table}[t]
        \centering
        \begin{tabular}{lllc}
            \toprule
            $(l,j)$ & $\mapsto$ & $(l',j')$ & $a_L$ \\
            \midrule
            $(l,l+s)$ & $\mapsto$ & $(l+1,l+s)$  & $\frac{1}{(2l+3)(2l+1)}$ \\
            $(l,l-s)$ & $\mapsto$ & $(l+1,l+s)$  & $\frac{l+1}{2l+1}$       \\
            $(l,l+s)$ & $\mapsto$ & $(l+1,l+3s)$ & $\frac{l+2}{2l+3}$       \\
            $(l,l-s)$ & $\mapsto$ & $(l-1,l-s)$  & $\frac{1}{(2l+1)(2l-1)}$ \\
            $(l,l+s)$ & $\mapsto$ & $(l-1,l-s)$  & $\frac{l}{2l+1}$         \\
            $(l,l-s)$ & $\mapsto$ & $(l-1,l-3s)$ & $\frac{l-1}{2l-1}$       \\
            \bottomrule
        \end{tabular}
        \caption{\textbf{Angular factors.} Table containing the angular factors
                 $a_l$ [see Eq.~\eqref{eq:OscillatorStrengthSimplified}] for
                 computing oscillator strengths. Throughout this table,
                 $s=1\slash{2}$.}
        \label{tab:AngularFactors}
    \end{table}
}



%%% START OF DOCUMENT
\begin{document}

    %%% TITLE PAGE
    \renewcommand{\pagemark}{}
    \maketitle%
    \newpage%

    %%% SET PLAIN PAGE STYLE
    \pagestyle{plain.scrheadings}%

    %%% TABLE OF CONTENTS

    % LINE FILLING
    \def\TOCLineFilling{\xleaders\hbox to.6em{\hfil{\tiny$\circ$\hfil}}\hfil}%

    % SET STYLES
    \DeclareTOCStyleEntry[linefill=\TOCLineFilling]{tocline}{section}%

    % PRINT
    \vspace*{2cm}
    \tableofcontents%
    \vfill
    {Typeset in \emph{Linux Libertine} with pdf-\LaTeXe \ in combination with
     \KOMAScript \ and \emph{latexmk}}
    \newpage

    %%% SET PAGE STYLE
    \pagestyle{scrheadings}%
    \setcounter{page}{1}%
    \renewcommand{\pagemark}{\thepage}%

    %%% SECTION
    \section{Introduction}
    \label{sec:introduction}

    Bound single-electron eigenfunctions of alkali-metal atoms and
    alkaline-earth metal ions are of relevance whenever properties of Rydberg
    atoms are of interest (e.g. in quantum computing and information, quantum
    simulation, Rydberg physics, quantum optics, \dots). Generally, because,
    except for the Hydrogen atom, the time-independent Schr\"odinger equation
    cannot be solved analytically, there is the need to obtain the eigenenergies
    (atomic spectrum) and eigenstates numerically. The library \alkcalc is
    intended for this purpose and information is given in the various provided
    \texttt{README} files. However, there is no discussion on the theory behind
    the code.

    This document's purpose is to define the Hamiltonian which \alkcalc is
    designed to diagonalise. Moreover, it formally defines the numerically
    computed quantities and hence serves as a reference for the user. In
    particular, the final section is a reference manual or the library
    functions.

    %%% SECTION
    \section{Theory}
    \label{sec:Theory}

    We start by considering a single atom or ion. Then we describe the \ac{fem}
    used to numerically solve the Schr\"odinger equation. Finally, we describe
    how we compute radial matrix elements, oscillator strengths, and
    state lifetimes.

    %%% SUBSECTION
    \subsection{Single-ion theory}
    \label{subsec:SingleIonTheory}

    The starting point is an alkali-metal atom or alkaline-earth-metal ion. We
    are interested in systems where all electrons are in closed shells except
    for one (e.g., $\mathrm{^1H}$, $\mathrm{^{38}Rb}$, $\mathrm{^{88}Sr^+}$).
    This allows a Kepler-problem description for the `effective nucleus' (the
    nucleus `screened' by the closed-shell electrons) and the valence electron.
    To be concrete, we study the Hamiltonian
    \begin{equation}
        H = \frac{\bm{p}^2}{2m}+V(\norm*{\bm{r}};\bm{L}^2,\bm{S}^2,\bm{J}^2)
    \,,
    \end{equation}
    where $m=Mm_\mathrm{e}\slash(M+m_\mathrm{e})$ is the reduced mass of the
    system, composed of the electron's mass $m_\mathrm{e}$ and the atom's or
    ion's mass $M$ without the valence electron. The centre-of-mass motion is
    decoupled, such that $\bm{r},\bm{p}$ are the relative canonical
    co-ordinates, i.e., $\bm{r}$ is the distance vector between the
    centre-of-mass position of the effective nucleus and the valence electron.
    Moreover, $\bm{L}$, $\bm{S}$, and $\bm{J}=\bm{L}+\bm{S}$ are the orbital
    angular momentum operator, the electron's spin operator, and the total
    angular momentum operator, respectively. The potential $V$ is not a simple
    Coulomb potential, but a generalised potential which includes different
    contributions (these arise when one starts with the full many-body electron
    Hamiltonian and incorporates relativistic and quantum electrodynamics
    effects). Following Refs.~\cite{aymar1996,marinescu1994,wilkinson2025} we
    write:
    \begin{equation}
        V(\norm*{\bm{r}};\bm{L}^2,\bm{S}^2,\bm{J}^2)
        =
        V_\mathrm{C}(\bm{r})
        +
        V_\mathrm{P}(\bm{r})
        +
        V_\mathrm{R}(\bm{r};\bm{L}^2,\bm{S}^2,\bm{J}^2)
    \,.
    \end{equation}
    The first contribution is
    \begin{equation}
        V_\mathrm{C}(r)
        =
        -\frac{e^2}{4\pii\varepsilon_0}\frac{Z_\mathrm{n}(r)}{r}
    \,, \
        Z_\mathrm{n}(r)
        =
        Z_\mathrm{c}+(Z-Z_\mathrm{c})\ee^{-k_1r}
        +
        k_2r\ee^{-k_3r}
        -
        k_4r^2\ee^{-k_3r}
    \,.
    \end{equation}
    This is a modified Coulomb potential, adjusted for small radii due to the
    effective nucleus. Here, $Z$ is the nuclear charge, $Z_\mathrm{c}$ the
    nuclear charge of the effective nucleus, and the $k_j$ ($j=1,2,3,4$) are
    empirically determined fitting parameters; see
    Refs.~\cite{aymar1996,marinescu1994,wilkinson2025}. The second contribution
    is
    \begin{equation}
        V_\mathrm{P}(r)
        =
        -\frac{e^2}{(4\pii\varepsilon_0)^2}
        \frac{\alpha_\mathrm{d}}{2r^4}
        \bigg[1-\exp(-\frac{r^6}{r^6_\mathrm{c}})\bigg]
    \,,
    \end{equation}
    which accounts for the polarisation of the screened nucleus due to the
    presence of the valence electron. Here, $r_\mathrm{c}$ is the effective
    nucleus' size and $\alpha_\mathrm{d}$ its polariseability. Finally,
    \begin{equation}\label{eq:vrIons}
        V_\mathrm{R}(r;\bm{L}^2,\bm{S}^2,\bm{J}^2)
        =
        \frac{1}{2m^2c^2}\frac{V'_\mathrm{NR}(r)}{rN(r)}\bm{L}\cdot\bm{S}
    \,, \
        N(r) = \bigg[ 1 - \frac{V_\mathrm{NR}(r)}{2mc^2} \bigg]^2
    \,, \
        V_\mathrm{NR}(r) = V_\mathrm{C}(r) + V_\mathrm{P}(r)
    \,,
    \end{equation}
    which describes (relativistic) coupling between the orbital angular momentum
    and the valence electrons' spin. In this expression, $V_\mathrm{NR}$ is the
    non-relativistic model potential and $c$ the speed of light. In
    Refs.~\cite{wilkinson2025,aymar1996} this potential is used for ions. We
    note that in Ref.~\cite{sibalic2017} the simplified version
    \begin{equation}\label{eq:vrAtoms}
        V_\mathrm{R}(r)
        =
        \frac{\alpha}{2r^3} \bm{L}\cdot\bm{S}
    \end{equation}
    is employed for atoms. Here, $\alpha$ is the fine-structure constant.
    However, we use the first, more precise, version for both ions and atoms.
    The spin-operator $\bm{S}$ in Eqs.~(\ref{eq:vrAtoms},\,\ref{eq:vrIons}) is
    conveniently expressed as
    \begin{equation}
        \bm{S} = \frac{\hbar}{2}\mqty[\sigma_x \\ \sigma_y \\ \sigma_z]
    \,,
    \end{equation}
    with the Pauli matrices
    \begin{equation}
        \sigma_x = \mqty[0&1 \\ 1&0]
    \,, \
        \sigma_y = \mqty[0&-\ii \\ \ii&0]
    \,, \
        \sigma_z = \mqty[1&0 \\ 0&-1]
    \,,
    \end{equation}
    which is possible because we deal with a single electron, i.e., a particle
    with spin $1\slash{2}$. Moreover, we rewrite the product $\bm{L}\cdot\bm{S}$
    in terms of the total angular momentum operator $\bm{J}=\bm{L}+\bm{S}$:
    \begin{equation}
        \bm{L}\cdot\bm{S}
        =
        \frac{1}{2}\big[\bm{J}^2-\bm{L}^2-\bm{S}^2\big]
    \,.
    \end{equation}
    For further simplification the square of the momentum operator is expressed
    in spherical co-ordinates,
    \begin{equation}
        \bm{p}^2
        =
        (-\ii\hbar\nabla)^2
        =
        -\hbar^2\nabla^2
        =
        -\frac{\hbar^2}{r^2}\pdv{r}(r^2\pdv{r})+\frac{\bm{L}^2}{r^2}
    \,.
    \end{equation}
    The second contribution is the standard centrifugal energy contribution,
    analogously to the Kepler problem. With these simplification we rewrite the
    Hamiltonian as
    \begin{equation}
        H
        =
        -\frac{\hbar^2}{2mr^2}\pdv{r}(r^2\pdv{r})
        +
        V_\mathrm{eff}(r;\bm{L}^2,\bm{S}^2,\bm{J}^2)
    \,, \
        V_\mathrm{eff}(r;\bm{L}^2,\bm{S}^2,\bm{J}^2)
        =
        \frac{\bm{L}^2}{2mr^2}+V(r;\bm{L}^2,\bm{S}^2,\bm{J}^2)
    \,.
    \end{equation}
    The full Hilbert space this Hamiltonian acts on is split into an angular
    momentum factor and a radial one. By construction, a basis of the angular
    momentum Hilbert space is given by the states $\ket*{l,s,j,m_j}$, which are
    eigenstates of $\bm{L}^2$, $\bm{S}^2$, $\bm{J}^2$, and $J_z$, i.e.,
    constitute the coupled basis, also called `fine-structure' basis. Note that
    it must be possible to construct an eigenbasis of $\bm{L}^2$, $\bm{S}^2$,
    $\bm{J}^2$, and $J_z$ simultaneously, because they commute pairwise. In
    particular, these states are constructed as linear combinations of the
    uncoupled basis states $Y_{l,m_l}\otimes\bm{\chi}_{s,m_s}$, employing
    Clebsch-Gordan coefficients. Here,
    $Y_{l,m_l}\colon[0,\pii]\times[0,2\pii]\to\mathbb{C}$ are spherical
    harmonics defined as (the so-called `Condon-Shortley phase' we include in
    the definition of the Legendre polynomials below, i.e., the factor
    $(-1)^{m_l}$ is \emph{not} missing here; see e.g. Ref.~\cite{friedrich2017})
    \begin{equation}\label{eq:sphericalHarmonics}
        Y_{l,m_l}(\vartheta,\varphi)
        =
        \sqrt{\frac{2l+1}{4\pii}\frac{(l-m_l)!}{(l+m_l)!}}
        P^{m_l}_l(\cos(\vartheta))\,\ee^{\ii m_l \varphi}
    \,,
    \end{equation}
    where the associated Legendre polynomials
    $P^{m_l}_l\colon[-1,1]\to\mathbb{R}$ are defined for $l=0,1,2,\dots$ and
    $m_l=0,1,\dots,l-1,l$, via the recursion relation
    \begin{equation}
        (l-m_l) P^{m_l}_l(x)
        =
        x (2l-1) P^{m_l}_{l-1}(x)
        -
        (l+m_l-1) P^{m_l}_{l-2}(x)
    \,,
    \end{equation}
    with starting point
    \begin{equation}
        P^{m_l}_{m_l}(x)
        =
        (-1)^{m_l}\frac{(2m_l)!}{2^{m_l}m_l!} (1-x^2)^{m_l\slash{2}}
    \,,
    \end{equation}
    in combination with $P^{m_l}_l(x)=0$ for $l<m_j$. Negative values of $m_l$
    are treated with the relation
    \begin{equation}
        P^{-m_l}_l = (-1)^{m_l} \frac{(l-m_l)!}{(l+m_l)!}P^{m_l}_l
    \,.
    \end{equation}
    Moreover, $\bm{\chi}_{s,m_s}$ are two-spinors defined as
    \begin{equation}\label{eq:twoSpinsors}
        \bm{\chi}_{\frac{1}{2},\frac{1}{2}}=\mqty[1 \\ 0]
    \,, \
        \bm{\chi}_{\frac{1}{2},-\frac{1}{2}}=\mqty[0 \\ 1]
    \,.
    \end{equation}
    Explicitly (for the case of interest $s=1\slash{2}$), we find the following
    expressions for the coupled basis states $\bm{\Phi}_{l,s,j,m_j}$:
    \begin{equation}\label{eq:coupledBasis}
        \bm{\Phi}_{l,s,j,m_j}
        =
        \langle Y_{l,m_j-s}\otimes\bm{\chi}_{s,s},\bm{\Phi}_{l,s,j,m_j} \rangle
        Y_{l,m_j-s}\otimes\bm{\chi}_{s,s}
        +
        \langle Y_{l,m_j+s}\otimes\bm{\chi}_{s,-s},\bm{\Phi}_{l,s,j,m_j} \rangle
        Y_{l,m_j+s}\otimes\bm{\chi}_{s,-s}
    \,,
    \end{equation}
    where the scalar products are (possibly vanishing) Clebsch-Gordan
    coefficients. We further find that these states transform according to the
    irreducible representations of the group $\mathrm{SU}(2)$ under the action
    of the $\mathrm{SU}(2)$-representation
    $\rho(\bm{\alpha})=\exp(\ii\bm{\alpha}\cdot\bm{J})$, where
    $\bm{\alpha}\in\mathbb{R}^3$ is a vector of length
    $\norm*{\bm{\alpha}}\in[0,2\pii)$. Moreover, the angular momentum
    Hamiltonian is invariant under the action of $\rho$, such that it is
    automatically block diagonal. The irreducible representations of
    $\mathrm{SU}(2)$ are $j(j+1)$-dimensional, with
    $j=0,1\slash{2},1,3\slash{2},2,\dots$, such that the blocks have dimensions
    $\mu_j{j}(j+1)$, where $\mu_j$ is the multiplicity of the irreducible
    representation $j$ in the direct sum decomposition of $\rho$. The
    eigenvalues associated to the blocks are operators on the radial Hilbert
    space. The multiplicities are $\mu_{j}=2$ and $\mu_0=0$, because for each
    fixed $j$ there are two possibilities: $l=j-1\slash{2}$ and
    $l=j+1\slash{2}$. Explicitly, the Hamiltonian is written as
    \begin{equation}
        \bigoplus^j_{m_j=-j}
        \langle \bm{\Phi}_{l,s,j,m_j}, H \bm{\Phi}_{l,s,j,m_j} \rangle
        \equiv
        \bigoplus^j_{m_j=-j} H_{l,s,j}
    \,,
    \end{equation}
    where (with the arbitrary choice $m_j=j$) we defined
    \begin{equation}
        H_{l,s,j}
        =
        \langle \bm{\Phi}_{l,s,j,j}, H \bm{\Phi}_{l,s,j,j} \rangle
        =
        -\frac{\hbar^2}{2mr^2}\pdv{r}(r^2\pdv{r})
        +
        V_\mathrm{eff}(r;\hbar^2l[l+1],\hbar^2s[s+1],\hbar^2j[j+1])
    \,.
    \end{equation}
    We note that the fact that the matrix element does not depend on $m_j$ is a
    consequence of the $\mathrm{SU}(2)$-symmetry of the Hamiltonian. Therefore,
    exactly this is the reason for the energies being degenerate in the total
    magnetic angular momentum quantum number $m_j$\footnote{Generally, there is
    no reason that the blocks $H_{l,s,j}$ are separate blocks for
    $l=j+1\slash{2}$ and $l=j-1\slash{2}$, because they correspond to the same
    irreducible representation $j$. However, the states we choose are already
    chosen in a way such that the blocks are already separated. In this case we
    know how we should choose the states already from the fact that the
    Hamiltonian explicitly depends on $\bm{L}$: if it would not be so obvious,
    we would not immediately know what the multiplicity index $l$ (and $s$)
    should physically be to sort the eigenvalues into separate blocks.} The
    symmetry consideration up until now simplified the problem to the block
    associated to the irreducible representation of $\mathrm{SU}(2)$. What is
    left to solve are the eigenproblems
    \begin{equation}\label{eq:radialEigenproblem}
        H_{l,s,j} R_{n,l,s,j} = E_{n,l,s,j} R_{n,l,s,j}
    \,,
    \end{equation}
    where $R_{n,l,s}$ is the $n$-th (radial) eigenstate of $H_{l,s,j}$.
    Therefore, if we manage to obtain $R_{n,l,s,j}$ the full spectrum is known.
    Of course, resonances (states associated to the channel
    $\bm{\Phi}_{l,s,j,m_j}$ which are not bound) exist, but in the following we
    will focus on bound states.

    %%% SUBSECTION
    \subsection{Numerical method}
    \label{subsec:NumericalMathod}

    The radial eigenstates, generally, cannot be obtained analytically.
    \alkcalc provides a platform to solve for these states numerically. To make
    the problem addressable by a computer, we express the eigenproblem in terms
    of dimensionless quantities. To this end we express the Hamiltonians
    $H_{l,s,j}$ in units of Hartree,
    $E_\mathrm{H}=1\slash{(4\pii\varepsilon_0a_\mathrm{B})}$, where
    $a_\mathrm{B}=\hbar\slash{(m_\mathrm{e}c\alpha)}\approx52.9\,\mathrm{pm}$ is
    Bohr's radius. The radial eigenfunctions we parametrise as
    \begin{equation}\label{eq:fRconversion}
        R_{n,l,s,j}(r)
        =
        \frac{1}{\sqrt{\aB^3}}\frac{f_{n,l,s,j}(r\slash\aB)}{r\slash\aB}
        \,.
    \end{equation}
    where $f_{n,l,s,j}$ is dimensionless. The eigenproblem in
    Eq.~\eqref{eq:radialEigenproblem} in the new eigenfunctions takes the form
    \begin{equation}\label{eq:eigenproblem}
        \bigg[
        -
        \frac{1}{2}
        \dv[2]{t}
        +\frac{l(l+1)}{2t^2}
        +C\frac{V(\aB{t};\hbar^2l[l+1],\hbar^2s[s+1],\hbar^2j[j+1])}{\EH}
        \bigg]
        f_{n,l,s,j}(t)
        =
        \bar{E}_{n,l,s,j}
        f_{n,l,s,j}(t)
    \,,
    \end{equation}
    where we defined
    \begin{equation}
        C = \frac{1}{1+\me\slash{M}}
    \end{equation}
    and the eigenenergies may be extracted from the relation
    \begin{equation}
        \frac{E_{l,s,j,n}}{E_\mathrm{H}}
        =
        \frac{\bar{E}_{n,l,s,j}}{C}
    \,.
    \end{equation}
    Rewriting the normalisation condition in terms of the eigenfunctions
    $f_{n,l,s,j}$ leads to
    \begin{equation}\label{eq:normalisation}
        1
        =
        \int^\infty_0\dd{r}r^2R^2_{n,l,s,j}(r)
        =
        \int^\infty_0 (\aB\dd{t})(\aB t)^2 \frac{1}{\aB^3}
        \frac{f^2_{n,l,s,j}(t)}{t^2}
        =
        \int^\infty_0\dd{t}f^2_{n,l,s,j}(t)
    \,.
    \end{equation}
    Here we assumed that the (radial) eigenfunctions are real-valued, which is
    justified as the Hamiltonian is not only an hermitian but also a symmetric
    operator. In conclusion, computing the normalised eigenfunctions
    $f_{n,l,s,j}$ yields the (automatically normalised) radial eigenstates by
    employing Eq.~\eqref{eq:fRconversion} for the conversion.

    To solve the eigenproblem in Eq.~\eqref{eq:eigenproblem} numerically, we
    employ a \ac{fem} instead of the conceptually (and programmatically) more
    simple \ac{fdm} for, mainly\footnote{It also helps to save some disk space,
    because one must store much less data to represent the eigenfunctions
    $f_{n,l,s,j}$. Typically, disk space is not a relevant factor these days,
    but we see no need to be unnecessarily wasteful with it either.}, one
    reason: With a \ac{fem} we are completely free in choosing the mesh (the
    position of the discretisation points) for the radial domain, while keeping
    the associated finite-dimensional matrix problem real and symmetric. This
    is important because the problem we face presents itself with drastically
    different energy scales, introduced by the diverging potential: the
    eigenfunctions $f_{n,l,s,j}$ oscillate quickly for small arguments and
    slowly for large ones, where the valence electron is almost free. Thus, a
    mesh tailored to this behaviour leads to a significant reduction of
    computational resources (the fast oscillation can be efficiently resolved
    without introducing `too many' points in the slowly oscillating regime).

    The first step in setting-up the \ac{fem} is to agree on a finite closed
    interval $I=[0, t_\mathrm{tmax}]$, where $t_\mathrm{max}>0$ sets the maximal
    value of the domain we consider. The value $t_\mathrm{max}$ must be chosen
    large enough to ensure that the boundary condition
    $f_{n,l,s,j}(t\geq{t}_\mathrm{max})=0$ is fulfilled. This equal sign is
    actually never true, because the radial eigenstates asymptotically approach
    zero. However, if we reach zero up to machine precision, it is numerically
    true. Of course, one must ensure that the finite set of solutions that shall
    be computed has support within this interval. The mesh we define as a family
    $(t_k)_{k=1,\dots,N-1}$ in $I$ with the properties $t_0=0$,
    $t_{N-1}=t_\mathrm{max}$, and $t_0<t_1<\dots<t_{N-2}<t_{N-1}$, where we say
    $N\geq{3}$. Of course, typical values are $N\,{\sim}\,10^5\text{-}10^7$.
    Next, we define the derived family of step sizes $(h_k)_{k=1,\dots,N-1}$,
    $h_k=t_k-t_{k-1}$. The elements $(\phi_k)_{k=0,\dots,N-1}$,
    $\phi_k\colon{I}\to[0,1]$, we define as
    \begin{equation}
        \phi_k(t)
        =
        \begin{cases}
            (t-t_{k-1})\slash(t_k-t_{k-1})\,, \ & t\in(t_{k-1},t_k)
        \\
            (t_{k+1}-t)\slash(t_{k+1}-t_k)\,, \ & t\in(t_k,t_{k+1})
        \\  0                             \,, \ & \text{else}
        \end{cases}
    \,.
    \end{equation}
    for $k=1,\dots,N-2$, and for $k=0$ and $k=N-1$ the Dirichlet boundary
    conditions $f_{n,l,s,j}(0)=f_{n,l,s,j}(t_\mathrm{max})$ lead to the
    canonical choice $\phi_0=\phi_{N-1}\equiv0$. To obtain a matrix problem
    which approximates Eq.~\eqref{eq:eigenproblem} we compute projections of it
    onto the elements with the real scalar product (`dummy' functions
    $\alpha,\beta\colon{I}\to\mathbb{R}$)
    \begin{equation}
        B(\alpha,\beta)
        =
        \int_I \dd{t} \alpha(t) \beta(t)
    \,.
    \end{equation}
    To simplify the notation we fix the quantum numbers $n$, $l$, $s$, and $j$
    and say $f\equiv{f_{n,l,s,j}}$. Moreover, we collect all potential
    contributions into the new potential $v$, defined as
    \begin{equation}
        v(t)
        =
        \frac{l(l+1)}{2t^2}
        +C\frac{V(\aB{t};\hbar^2l[l+1],\hbar^2s[s+1],\hbar^2j[j+1])}{\EH}
    \,.
    \end{equation}%
    %
    Finally, we say $\lambda=\bar{E}_{n,l,s,j}$, such that
    Eq.~\eqref{eq:eigenproblem} simply becomes
    \begin{equation}
        -\frac{1}{2}f''(t) + v(t) f(t) = \lambda f(t)
    \,.
    \end{equation}
    Before we project this equation onto the elements, we approximately expand
    $f$ in them:
    \begin{equation}\label{eq:expansionOfState}
        f = \sum^{N-2}_{k=1} f_k \phi_k
    \,.
    \end{equation}
    This expansion automatically fulfils the Dirichlet boundary conditions
    $f(0)=f(t_\mathrm{max})=0$. Projection, i.e., applying the map
    $B(\phi_k,\cdot)$ for all $k=1,\dots,N-2$, yields a matrix eigenproblem
    who's solutions converge to solutions of Eq.~\eqref{eq:eigenproblem} as $N$
    tends to infinity:
    \begin{equation}\label{eq:matrixEigenproblem}
        (K+V) \bm{f} = \lambda M \bm{f}
    \,.
    \end{equation}
    In this expression $M_{kk'}=B(\phi_k,\phi_{k'})$ is the so-called mass
    matrix, $K_{kk'}=B(\phi'_k,\phi'_{k'})\slash{2}$ the so-called stiffness
    matrix\footnote{Note that because the elements obey the Dirichlet boundary
    conditions $\phi_k(0)=\phi_k(t_\mathrm{max})=0$, $k=1,\dots,N-2$,
    integrating by parts allows to write
    $B(\phi''_k,\phi_{k'})=-B(\phi'_k,\phi'_{k'})$.}, and
    $V_{kk'}=B(\phi_k,v\phi_{k'})$ the potential energy contribution. Each
    matrix is $(N-2)$-dimensional, since $k,k'=1,\dots,N-2$, and $M$ and $K$ can
    be computed analytically; in particular
    \begin{equation}
        M
        =
        \frac{1}{6}
        \mqty[2(h_1+h_2) &  h_2 & & & \\
              h_2 & 2(h_2+h_3) & h_3 & & \\
              & h_3 & 2(h_3+h_4) & h_4 & \\
              & & \ddots & \ddots & h_{N-2} \\
              & & & h_{N-2} & 2(h_{N-2}+h_{N-1})]
    \end{equation}
    and
    \begin{equation}
        K = \frac{1}{2}
            \mqty[\frac{1}{h_1}+\frac{1}{h_2} & -\frac{1}{h_2} & \\
                  -\frac{1}{h_2} & \frac{1}{h_2}+\frac{1}{h_3} & -\frac{1}{h_3}
                  & \\
                  & -\frac{1}{h_3} & \frac{1}{h_3}+\frac{1}{h_4} &
                  -\frac{1}{h_4} & \\
                  & & \ddots & \ddots & -\frac{1}{h_{N-2}} \\
                  & & & -\frac{1}{h_{N-2}} & \frac{1}{h_{N-2}}+\frac{1}{h_{N-1}}
                  ]
    \,.
    \end{equation}
    Note that both of these matrices are symmetric. The potential matrix,
    generally, must be computed numerically. However, here we choose the simple
    approximation
    \begin{equation}
        B(\phi_k,v\phi_{k'})
        \approx
        \frac{v(t_k)+v(t_{k'})}{2}
        M_{k{k'}}
    \,,
    \end{equation}
    which is justified if the step size is small enough. Note that this
    particular choice of approximations keeps the matrix symmetric, such that
    $K+V$ is symmetric. Furthermore, $M$ represents a scalar product and is thus
    positive definite. This means that the matrix eigenproblem in
    Eq.~\eqref{eq:matrixEigenproblem} can be efficiently solved with the Lanczos
    algorithm.

    Next, we express the normalisation condition in
    Eq.~\eqref{eq:normalisation} in terms of our elements and the vector
    $\bm{f}$. To this end we calculate
    \begin{equation}
        \int^\infty_0\dd{t}f^2(t)
        \approx
        \int_I  \dd{t} f^2(t)
        =
        \sum^{N-2}_{k=1}\sum^{N-2}_{k'=1}f_kf_{k'} B(\phi_k,\phi_{k'})
        =
        \bm{f}^\mathrm{T}M\bm{f}
        =
        \norm*{\bm{f}}^2_M
        \overset{!}{=}
        1
    \,,
    \end{equation}
    i.e., the length of $\bm{f}$ must be unity in the by the scalar product $M$
    induced norm $\norm*{\,\cdot\,}_M$.

    %%% SUBSECTION
    \subsection{Radial matrix elements}
    \label{subsec:RadialMatrixElements}

    To compute radial matrix elements, we first calculate, for $p\in\mathbb{R}$
    such that the integral exists,
    \begin{align}\label{eq:radialMatrixElement}
        \begin{split}
            \frac{r^{n,l,s,j;n',l',s',j'}_p}{\aB^p}
            &\equiv
            \frac{1}{\aB^p} \langle R_{n,l,s,j},
            \mathrm{id}^p R_{n',l',s',j'} \rangle
            =
            \int^\infty_0 \dd{r} r^2
            R_{n,l,s,j}(r)
            \frac{r^p}{\aB^p}
            R_{n',l',s',j'}(r)
        \\
            &=
            \int^\infty_0 (\dd{t} \aB) (\aB t)^2
            f_{n,l,s,j}(t)
            \frac{(\aB t)^p}{\aB^p} \frac{1}{\aB^3}
            \frac{f_{n',l',s',j'}(t)}{t^2}
            =
            B(f_{n,l,s,j}, \mathrm{id}^p f_{n',l',s',j'})
        \,.
        \end{split}
    \end{align}
    Here, `$\mathrm{id}$' denotes the identity on the non-negative real numbers.
    Next, we approximate this result with solutions to the matrix eigenproblem
    in Eq.~\eqref{eq:matrixEigenproblem}, i.e., we represent $f_{n,l,s,j}$ and
    $f_{n',l',s',j'}$ approximately with the vectors $\bm{f}_{n,l,s,j}$ and
    $\bm{f}_{n',l',s',j'}$, respectively [see Eq.~\eqref{eq:expansionOfState}].
    This allows to write
    \begin{equation}
        B(f_{n,l,s,j}, \mathrm{id}^p f_{n',l',s',j'})
        \approx
        \bm{f}_{n,l,s,j}^\mathrm{T} R_p \bm{f}_{n',l',s',j'}
    \,,
    \end{equation}
    where the components of the real symmetric matrix $R_p$ are defined as
    \begin{equation}
        (R_p)_{k,k'}
        =
        B(\phi_k,\mathrm{id}^p\phi_{k'})
        =
        \begin{cases}
            B(\phi_k,\mathrm{id}^p\phi_{k+1}),
            & k'=k+1 \ \text{or} \ k=k'+1 \ (0<k,k'<N-2)
        \\
            B(\phi_k,\mathrm{id}^p\phi_{k}),
            & k'=k \ (0<k,k'\leq{N-2})
        \\
        0, & \text{else}
        \end{cases}
    \,.
    \end{equation}
    To numerically compute the remaining integrals we use a two-point
    Gau\ss{i}an quadrature. This method approximates the integral as
    \cite{burden2011}:
    \begin{equation}
        \int^b_a \dd{t} \alpha(t)
        \approx
        \frac{b-a}{2}
        \bigg[
            \alpha\bigg(\frac{a+b}{2}-\frac{b-a}{2\sqrt{3}}\bigg)
            +
            \alpha\bigg(\frac{a+b}{2}+\frac{b-a}{2\sqrt{3}}\bigg)
        \bigg]
    \,,
    \end{equation}
    for any sufficiently smooth function $\alpha$ and $a<b$. Note that this
    result is exact if $\alpha$ is a polynomial of degree three, which is the
    case for radial dipole matrix elements. This method is straight-forward to
    implement in a computer program, such that the problem of calculating matrix
    elements is effectively reduced to computing matrix products. However, note
    that for too large $p$ there can be numerical overflow or underflow when
    computing $t^p$. But, for 64-bit floating point arithmetic, $p=\pm10$ is
    easily possible (if the integral exists), such that, typically, all
    practical cases are included. Explicitly, we could verify that for Hydrogen
    the range $p=-2,\dots,10$ yields reasonably good results. Finally, we remark
    that one can improve on the precision by employing higher-order Gau\ss{i}an
    quadrature algorithms as described in \cite{burden2011}.

    %%% SUBSECTION
    \subsection{Oscillator strengths}
    \label{subsec:OscillatorStrengths}

    The so-called `oscillator strength', for a transition from an initial state
    $\bm{\Psi}_\mathrm{i}\equiv\bm{\Psi}_{n',l',s,j',m'_j}$ to a final state
    $\bm{\Psi}_\mathrm{f}\equiv\bm{\Psi}_{n,l,s,j,m_j}$ is
    \cite{friedrich2017}
    \begin{equation}
        f_{\mathrm{i}\to\mathrm{f}}
        =
        \sum^{j'}_{m'_j=-j'}
        \frac{2}{3} \bar{E}_\mathrm{f,i}
        \abs*{\langle\bm{\Psi}_\mathrm{f},
                     \bm{r}\slash\aB\bm{\Psi}_\mathrm{i}\rangle}^2
    \,,
    \end{equation}
    with the energy difference
    $\EH\bar{E}_\mathrm{f,i}=E_\mathrm{f}-E_\mathrm{i}$. Note that we sum over
    the all final $m'_j$, because we are interested in transitions between the
    fine-structure states without resolving the Zeeman sub-levels. The oscillator
    strength is a measure for how strongly the induced Hertzian (electric)
    dipole oscillates: the radiation power is large for spatially large dipoles
    and fast oscillation frequencies, captured by the matrix element and energy
    difference, respectively. In principle, we already have all the machinery
    to numerically compute this quantity, however, one can reduce the expression
    to a much simpler form by using angular momentum algebra, in particular,
    relations from \cite{varshalovic1988}. To simplify the expression, we first
    use (Condon-Shortley phase)
    \begin{equation}
        r_x
        =
        \sqrt{\frac{2\pii}{3}} r (Y_{1,-1}-Y_{1,1})
    \,, \
        r_y
        =
        \ii \sqrt{\frac{2\pii}{3}} r (Y_{1,-1}+Y_{1,1})
    \,, \
        r_z
        =
        \sqrt{\frac{4\pii}{3}} r Y_{1,0}
    \,.
    \end{equation}
    With this, the oscillator strength becomes
    \begin{equation}
        f_{\mathrm{i}\to\mathrm{f}}
        =
        \frac{4\pii}{3} \frac{2}{3} \bar{E}_\mathrm{f,i}
        \langle{r}R_\mathrm{f},r^2\slash\aB{r}R_\mathrm{i}\rangle^2
        \sum^{j'}_{m'_j=-j'}
        \sum^1_{q=-1}
        \abs*{\langle\bm{\Phi}_\mathrm{f},Y_{1,q}\bm{\Phi}_\mathrm{i}\rangle}^2
    \end{equation}
    Employing relations from \cite[p.~236,244,260,445,496]{varshalovic1988} one
    can show that the remaining matrix element is
    \begin{equation}
        \langle \bm{\Phi}_\mathrm{f}, Y_{1,q}\bm{\Phi}_\mathrm{i} \rangle
        =
        (-1)^{j'+j+1+l'+s-m'_j}
        \sqrt{\frac{3}{4\pii}}
        \mqty(j' & 1 & j \\ -m'_j & q & m_j)
        \sqrt{(2j+1)(2j'+1)(2l+1)}
        \bigg\{\mqty{l & s & j \\ j' & 1 & l'}\bigg\}
        C^{l,0,1,0}_{l',0}
    \,,
    \end{equation}
    expressed in terms of a Wigner $3jm$-symbol, a Wigner $6j$-symbol, and a
    Clebsch-Gordan coefficient. With symmetry relations and a sum rule from
    \cite[p.~236,259]{varshalovic1988} we can compute the sum over $q$ and
    $m'_j$ in the expression for the oscillator strength, which yields
    \begin{equation}\label{eq:OscillatorStrengthSimplified}
        \sum_{m'_j} f_{\mathrm{i}\to\mathrm{f}}
        =
        \frac{2}{3} \bar{E}_\mathrm{f,i}
        \abs*{\langle r R_{n',l',s,j'}, r^2\slash\aB R_{n,l,s,j} \rangle}^2
        a_l
    \,, \
        a_l
        =
        l_\mathrm{max} (2j'+1)
        \bigg\{\mqty{l & s & j \\ j' & 1 & l'}\bigg\}^2
    \,,
    \end{equation}
    where $l_\mathrm{max}=\mathrm{max}(l,l')$. There are only six transitions
    for which $a_l$ is non-zero. Using relations from
    \cite[p.~298,300]{varshalovic1988} we arrive at the results shown in
    Tab.~\ref{tab:AngularFactors}. With this is is now straight-forward to
    compute the oscillator strengths numerically with a computer program. As a
    final remark, note that the oscillator strength does not depend on $m_j$ any
    more, i.e., by summing over final states, we find that the oscillators for
    different $m_j$ all have the same strength.

    %%% TABLE
    \TabAngularFactors

    %%% SUBSECTION
    \subsection{State lifetimes}
    \label{subsec:StateLifetimes}

    To compute the lifetime of a fine-structure state, we need to calculate the
    so-called `Einstein-A' coefficient \cite{einstein1917,friedrich2017}. It is
    a rate, that describes the number of spontaneous emissions per unit time at
    zero temperature. The total decay rate, at zero temperature, is given by the
    sum over all Einstein-A coefficients. In the dipole approximation this sum
    is given by the expression \cite{friedrich2017}:
    \begin{equation}\label{eq:DecayRateGeneralForm}
        A_\mathrm{i}
        =
        \sum_\mathrm{f}
        A_{\mathrm{i}\to\mathrm{f}}
    \,, \
        A_{\mathrm{i}\to\mathrm{f}}
        =
        \frac{\EH}{\hbar} \frac{4}{3} \alpha^2 \bar{E}^3_\mathrm{i,f}
        \sum_{u=x,y,z}
        \abs*{\langle\bm{\Psi}_\mathrm{f},r_u\bm{\Psi}_\mathrm{i}\rangle}^2
    \,.
    \end{equation}
    In terms of the oscillator strength we obtain
    \begin{equation}
        A_\mathrm{i}
        =
        \frac{2\alpha^3\EH}{\hbar}
        \sum_{n',l',j'}
        E^2_\mathrm{i,f} f_{\mathrm{f}\to\mathrm{i}}
    \,.
    \end{equation}
    This is the rate of spontaneous emission which only appears for transitions
    where $E_\mathrm{i}>E_\mathrm{f}$. Note that, implicitly, this restricts the
    quantum numbers the sum is running over. Further, this is the correct decay
    rate only for zero temperature case. At finite temperature, spontaneous
    transitions increasing the energy may be driven as well, if thermal
    fluctuations are sufficiently strong. The background spectral density is
    given by Plank's law, which yields for the average photon number at
    frequency $\omega$ and temperature $T$:
    \begin{equation}
        n(\omega,T)
        =
        \frac{1}{\exp(\hbar\omega\slash\kBT)-1}
    \,,
    \end{equation}
    where $k_\mathrm{B}$ is Boltzmann's constant. From Einstein's rate equations
    \cite{einstein1917} it follows that the total rate for decay out of the
    initial state is
    \begin{equation}
        \frac{\hbar}{\EH}\Gamma
        =
        2\alpha^3
        \sum_{n',l',j'}
        \bar{E}^2_\mathrm{i,f} f_{\mathrm{f}\to\mathrm{i}}
        (1+n(E_\mathrm{i,f}\slash\hbar,T))
        +
        \sum_{n',l',j'}
        \bar{E}^2_\mathrm{f,i} f_{\mathrm{i}\to\mathrm{f}}
        (1+n(E_\mathrm{f,i}\slash\hbar,T))
    \,,
    \end{equation}
    where each sum only runs over quantum numbers such that the associated term
    is positive. That is, the first sum is only over final states with lower,
    and the second one only over final states with larger, energy than the
    initial state. With this formula, the lifetime of a fine-structure state is
    $\tau=\Gamma^{-1}$.

    \clearpage

    %%% SECTION
    \section{Manual}
    \label{sec:Manual}

    Reference manual for the \alkcalc library functions. Prototypes for the
    following functions are defined in the header `interface.d/alkcalc.h'.

    \sffamily%
    \pagestyle{plain}%
    \vspace*{12mm}\par%

    \begin{sdocfun}
    {
        alkcalc\tus{Enlsj}
    }{
        Eigenenergies in units of Hartree. The requested eigenenergy is read
        from the corresponding data file in the directory `data.d' and returned.
    }
        \sdocarg{char *species}{
            String to specify atom/ion species, e.g., \texttt{1H} for Hydrogen,
            or \texttt{88SR+} for the $\mathrm{^{88}Sr^+}$ ion.
        }
        \sdocarg{int32\tus{t} n}{
            Principal quantum number.
        }
        \sdocarg{int32\tus{t} l}{
            Orbital angular momentum quantum number $l=1,2,\dots,n-1$.
        }
        \sdocarg{double s}{
            Electron spin quantum number. Not an argument here, since
            $s=1\slash{2}$ is always assumed.
        }
        \sdocarg{double j}{
            Total angular momentum quantum number
            $j=\abs*{l-1\slash{2}},l+1\slash{2}$. Here, the following convention
            is applied: For $j\in[k,k+1)$, $k=0,1,2,\dots$, the library
            internally uses the (half-integer) value $(2k+1)\slash{2}$. For
            instance, this means that for $j=0.49999$ and $j=0.50001$ the
            software internally uses $j=1\slash{2}$ in both cases.
        }
        \sdocret{double E}{
            $E$ is the eigenenergy specified by the quantum numbers $n$, $l$,
            $s=1\slash{2}$, and $j$.
        }
    \end{sdocfun}%
    %
    \begin{sdocfun}
    {
        alkcalc\tus{fnlsj}
    }{
        Numerical equivalent to $f_{n,l,s,j}$ in Eq.~\eqref{eq:fRconversion}.
        The data is read from the file containing the requested state. This data
        is stored at a user specified location. The location where \alkcalc will
        search for the data can be set in the file `interface.d/alkcalc.h'. The
        result is owned by the caller and should be freed with
        \texttt{alkcalc\tus{state}\tus{free}}.
    }
        \sdocarg{char result}{
            Character specifying if the discretisation data \texttt{t} and
            \texttt{h} is returned (see RETURN down below). If the argument
            \texttt{result} is equal to \texttt{f} (full), the full result is
            returned, and if \texttt{result} is equal to \texttt{p} (partial),
            \texttt{t=h=NULL}.
        }
        \sdocarg{char *species}{
            String to specify atom/ion species, e.g., \texttt{1H} for Hydrogen,
            or \texttt{88SR+} for the $\mathrm{^{88}Sr^+}$ ion.
        }
        \sdocarg{int32\tus{t} n}{
            Principal quantum number.
        }
        \sdocarg{int32\tus{t} l}{
            Orbital angular momentum quantum number $l=1,2,\dots,n-1$.
        }
        \sdocarg{double s}{
            Electron spin quantum number. Not an argument here, since
            $s=1\slash{2}$ is always assumed.
        }
        \sdocarg{double j}{
            Total angular momentum quantum number
            $j=\abs*{l-1\slash{2}},l+1\slash{2}$. Here, the following convention
            is applied: For $j\in[k,k+1)$, $k=0,1,2,\dots$, the library
            internally uses the (half-integer) value $(2k+1)\slash{2}$. For
            instance, this means that for $j=0.49999$ and $j=0.50001$ the
            software internally uses $j=1\slash{2}$ in both cases.
        }
        \sdocret{alkcalc\tus{state} *f}{
            \texttt{f} contains eight members: \texttt{N}, \texttt{dim},
            \texttt{t}, \texttt{h}, and \texttt{fnlsj}.\par%
            \texttt{int32\tus{t} N}: Number of discretisation points $N$ in the
            mesh, i.e., if $N=10$, the interval $[0,t_\mathrm{max}]$ is split
            into $N-1=10-1=9$ intervals.\par%
            \texttt{int32\tus{t} n}: Principal quantum number.\par%
            \texttt{int32\tus{t} l}: Orbital angular momentum quantum
            number.\par%
            \texttt{double j}: Total angular momentum quantum number.\par%
            \texttt{int32\tus{t} dim}: Dimension of eigenproblem, i.e., $N-2$,
            which is the number of elements appearing in the expansion in
            Eq.~\eqref{eq:expansionOfState}.\par%
            \texttt{double *t}: Array of length $N$, containing the mesh
            $(t_k)_{k=0,\dots,N-1}$.\par%
            \texttt{double *h}: Array of length $N-1$, containing the step
            sizes, $(h_k)_{k=1,\dots,N-1}$.\par%
            \texttt{double *fnlsj}: Array of length $N-2$ containing the
            requested eigenvector $\bm{f}$. The components of $\bm{f}$ are the
            coefficients $f_k$ ($k=1,\dots,N-2$) defining the expansion in
            Eq.~\eqref{eq:expansionOfState}.
        }
    \end{sdocfun}%
    %
    \begin{sdocfun}
    {
        alkcalc\tus{rp}
    }{
        Radial matrix element with arbitrary power $p\in\mathbb{R}$ [see
        Eq.~\eqref{eq:radialMatrixElement}]. For too large powers $\abs*{p}>10$,
        the there is the risk of numerical over or underflow of the 64-bit
        floating point arithmetic. Generally, $\mathrm{^1H}$ (Hydrogen atom) and
        $\mathrm{^2He^+}$ (Helium ion) can be used as a benchmark for atoms and
        ions, respectively. It is the user's responsibility to ensure that the
        integral associated to the desired matrix element exists. In case it
        does not, the returned value is wrong and no error is thrown.
    }
        \sdocarg{char *species}{
            String to specify atom/ion species, e.g., \texttt{1H} for Hydrogen,
            or \texttt{88SR+} for the $\mathrm{^{88}Sr^+}$ ion.
        }
        \sdocarg{int32\tus{t} nb}{
            Principal quantum number of bra.
        }
        \sdocarg{int32\tus{t} lb}{
            Orbital angular momentum quantum number $l=1,2,\dots,n-1$ of bra.
        }
        \sdocarg{double sb}{
            Electron spin quantum number of bra. Not an argument here, since
            $s=1\slash{2}$ is always assumed.
        }
        \sdocarg{double jb}{
            Total angular momentum quantum number of bra, $j$. Here, the
            following convention is applied: For $j\in[k,k+1)$, $k=0,1,2,\dots$,
            the library internally uses the (half-integer) value
            $(2k+1)\slash{2}$. For instance, this means that for $j=0.49999$ and
            $j=0.50001$ the software internally uses $j=1\slash{2}$ in both
            cases.
        }
        \sdocarg{double p}{
            Power of radius operator; cf. Eq.~\eqref{eq:radialMatrixElement}.
            For instance, $p=0$ yields the scalar product of the states. For
            $p=1$ and $p=2$, respectively, dipole and quadrupole matrix elements
            are computed. Generally, $p\in\mathbb{R}$.
        }
        \sdocarg{int32\tus{t} nk}{
            Principal quantum number of ket.
        }
        \sdocarg{int32\tus{t} lk}{
            Orbital angular momentum quantum number $l=1,2,\dots,n-1$ of ket.
        }
        \sdocarg{double sk}{
            Electron spin quantum number of ket. Not an argument here, since
            $s=1\slash{2}$ is always assumed.
        }
        \sdocarg{double jk}{
            Total angular momentum quantum number of ket, $j$. Here, the
            convention for \texttt{jb} (see above) is adopted.
        }
        \sdocret{double rp}{
            \texttt{rp} is the matrix element $r^{n,l,s,j;n',l',s',j'}_p$ in
            Eq.~\eqref{eq:radialMatrixElement} in units of $\aB^p$, where $\aB$
            is Bohr's radius. Quantum numbers without(with) a prime correspond
            to the quantum numbers of the bra(ket).
        }
    \end{sdocfun}%
    %
    \begin{sdocfun}
    {
        alkcalc\tus{cj1m1j2m2jmj}
    }{
        Clebsch-Gordan coefficients,
        \begin{equation}
            C^{j_1,m_1,j_2,m_2}_{j, m_j}
            =
            \braket*{j_1,m_1;j_2,m_2}{j_1,j_2,j,m_j}
        \,,
        \end{equation}%
        %
        where $j$ is the total angular momentum composed of the two individual
        ones $j_1$ and $j_2$. To fix the phase uniquely, we adopt the
        Condon-Shortley sign convention, that is, the Clebsch-Gordan
        coefficients with $m_j=j_1+j_2$ are positive. In particular, for the
        case $j_1=l$ and $j_2=s$,
        \begin{equation}\label{eq:definitionCGC}
            C^{l,m_l,s,m_s}_{j,m_j}
            =
            \langle
            Y_{l,m_l}\otimes\bm{\chi}_{s,m_s},
            \bm{\Phi}_{l,s,j,m_j}
            \rangle
        \,,
        \end{equation}%
        %
        with the states $Y_{l,m_l}\otimes\bm{\chi}_{s,m_s}$ of the uncoupled
        basis and the coupled basis states $\bm{\Phi}_{l,s,j,m}$; see
        Eqs.~(\ref{eq:sphericalHarmonics},\,\ref{eq:twoSpinsors}), and
        Eq.~\eqref{eq:coupledBasis}, respectively.  The return type (see RETURN
        down below) allows to extract the requested coefficient symbolically as
        well as numerically. The function returns zero in case the arguments do
        not make sense, e.g., if $j_1<m_1$; or if $j_1$ is half-integer and
        $m_1$ is integer; or if $j$ does not lay between (including bounds)
        $\abs*{j_1-j_2}$ and $\abs*{j_1+j_2}$. The arguments to this function
        are defined as double precision floating point numbers, in order to
        represent half-integer values. For numerical stability and definiteness
        the arguments (positive as well as negative) are read with the following
        convention: for $a=j_1,m_1,j_2,m_2,j,m_j$ the value
        $a=1\slash{2}\lfloor{2a+1\slash{2}}\rfloor$ is used, with
        $\lfloor{x}\rfloor=\mathrm{argmin}\{\abs*{x-z}\colon{z}\in\mathbb{Z},
        z\leq{x}\}$. This might seem somewhat convoluted, but simply means that
        for $-0.25\leq{a}<0.25$ we take $a=0$; for $0.25\leq{a}<0.75$ we take
        $a=1\slash{2}$; for $0.75\leq{a}<1.25$ we take $a=1$; and so on (in both
        positive and negative direction). This makes sure that the finite
        precision of double precision floating point numbers does not play a
        role; i.e., it does not matter if one puts $a=0.5$ or $a=0.500001$ or
        $a=0.499998$ as an argument --- $a=1\slash{2}$ is used internally in any
        case.
    }
        \sdocarg{double j1}{
            Angular momentum quantum number of first factor.
        }
        \sdocarg{double m1}{
            Magnetic quantum number of first factor.
        }
        \sdocarg{double j2}{
            Angular momentum quantum number of second factor.
        }
        \sdocarg{double m2}{
            Magnetic quantum number of second factor.
        }
        \sdocarg{double j}{
            Total angular momentum quantum number.
        }
        \sdocarg{double mj}{
            Magnetic quantum number of total angular momentum.
        }
        \sdocret{alkcalc\tus{cg} c}{
            \texttt{c} contains three members: \texttt{sign},
            \texttt{numerator}, and \texttt{denominator}.\par%
            \texttt{int8\tus{t} sign}: Sign, $\sigma$, of the Clebsch-Gordan
            coefficient.\par%
            \texttt{int64\tus{t} numerator}: Numerator, $\nu$, of the
            Clebsch-Gordan coefficient.\par%
            \texttt{int64\tus{t} denominator}: Denominator, $\mu$, of the
            Clebsch-Gordan coefficient.\par%
            The requested Clebsch-Gordan coefficient is given by
            $\sigma\sqrt{\nu\slash\mu}$. Euklid's algrithm is employed to reduce
            $\nu$ and $\mu$. This representation allows to obtain Clebsch-Gordan
            coefficients symbolically as well as numerically.
        }
    \end{sdocfun}%
    %
    \begin{sdocfun}
    {
        alkcalc\tus{YlmlXsms}
    }{
        Angular factor of the full eigenstate in the uncoupled basis, i.e., the
        two-component spinor $Y_{l,m_l}(\vartheta,\varphi)\bm{\chi}_{s,ms}$.
        The input quantum numbers are checked for inconsistencies and an error
        is thrown in case one is detected. The angles $\vartheta$ and $\varphi$
        have the ranges $[0,\pii]$ and $[0,2\pii]$, respectively. This condition
        is checked up to floating point machine precision (type \texttt{double})
        and an error will be thrown in case either condition is not obeyed.
    }
        \sdocarg{int32\tus{t} l}{
            Orbital angular momentum quantum number $l=1,2,\dots,n-1$.
        }
        \sdocarg{int32\tus{t} ml}{
            Magnetic quantum number of orbital angular momentum.
        }
        \sdocarg{double s}{
            Electron spin quantum number. Not an argument here, since
            $s=1\slash{2}$ is always assumed.
        }
        \sdocarg{double ms}{
            Magnetic quantum number of electron spin, $m_s$. For any $m_s<0$
            the value $m_s=-1\slash{2}$ is used internally, and for $m_s\geq0$
            the value $m_s=1\slash{2}$.
        }
        \sdocarg{double theta}{
            Polar angle (Zenitwinkel), $\vartheta$, in range $[0,\pii]$.
        }
        \sdocarg{double phi}{
            Azimuthal angle (Azimut), $\varphi$, in range $[0,2\pii]$.
        }
        \sdocret{alkcalc\tus{spinor} s}{
            \texttt{s} has two components and represents a two-component spinor:
            \texttt{u} is the spin-up ($m_s=1\slash{2}$) component and
            \texttt{d} is the spin-down ($m_s=-1\slash{2}$) component.
        }
    \end{sdocfun}%
    %
    \begin{sdocfun}
    {
        alkcalc\tus{Philsjmj}
    }{
        Angular factor of the full eigenstate in the coupled (fine-structure)
        basis, i.e., the two-component spinor
        $\bm{\Phi}_{l,s,j,mj}(\vartheta,\varphi)$ [see
        Eq.~\eqref{eq:coupledBasis}]. The input quantum numbers are checked for
        inconsistencies and an error is thrown in case one is detected. The
        angles $\vartheta$ and $\varphi$ have the ranges $[0,\pii]$ and
        $[0,2\pii]$, respectively. This condition is checked up to floating
        point machine precision (type \texttt{double}) and an error will be
        thrown in case either condition is not obeyed. To compute the resulting
        two-component spinor Clebsch-Gordon coefficients are needed, for which
        the sign convention detailed in Eq.~\eqref{eq:definitionCGC} is
        employed.
    }
        \sdocarg{int32\tus{t} l}{
            Orbital angular momentum quantum number $l=1,2,\dots,n-1$.
        }
        \sdocarg{double s}{
            Electron spin quantum number. Not an argument here, since
            $s=1\slash{2}$ is always assumed.
        }
        \sdocarg{double j}{
            Total angular momentum quantum number
            $j=\abs*{l-1\slash{2}},l+1\slash{2}$. Here, the following convention
            is applied: For $j\in[k,k+1)$, $k=0,1,2,\dots$, the library
            internally uses the (half-integer) value $(2k+1)\slash{2}$. For
            instance, this means that for $j=0.49999$ and $j=0.50001$ the
            software internally uses $j=1\slash{2}$ in both cases.
        }
        \sdocarg{double mj}{
            Magnetic quantum number of total angular momentum, $m_j$. Regarding
            the floating point representation of half-integers, the same
            convention as for the total angular momentum quantum number $j$ (see
            above) is employed.
        }
        \sdocarg{double theta}{
            Polar angle (Zenitwinkel), $\vartheta$, in range $[0,\pii]$.
        }
        \sdocarg{double phi}{
            Azimuthal angle (Azimut), $\varphi$, in range $[0,2\pii]$.
        }
        \sdocret{alkcalc\tus{spinor} s}{
            \texttt{s} has two components and represents a two-component spinor:
            \texttt{u} is the spin-up ($m_s=1\slash{2}$) component and
            \texttt{d} is the spin-down ($m_s=-1\slash{2}$) component.
        }
    \end{sdocfun}%
    %
    \begin{sdocfun}
    {
        alkcalc\tus{fitof}
    }{
        Oscillator strength, for the transition between an initial state
        $\bm{\Psi}_\mathrm{i}$ and a final state $\bm{\Psi}_\mathrm{f}$, as
        defined in Subsec.~\ref{subsec:OscillatorStrengths}, that is, it was
        already summed over the magnetic quantum number $m'_j$ of the final
        state.
    }
        \sdocarg{char *species}{
            String to specify atom/ion species, e.g., \texttt{1H} for Hydrogen,
            or \texttt{88SR+} for the $\mathrm{^{88}Sr^+}$ ion.
        }
        \sdocarg{int32\tus{t} ni}{
            Principal quantum number of the initial state
            $\bm{\Psi}_\mathrm{i}$.
        }
        \sdocarg{int32\tus{t} li}{
            Orbital angular momentum quantum number $l=1,2,\dots,n-1$ of the
            initial state $\bm{\Psi}_\mathrm{i}$.
        }
        \sdocarg{double si}{
            Electron spin quantum number of the initial state
            $\bm{\Psi}_\mathrm{i}$. Not an argument here, since $s=1\slash{2}$
            is always assumed.
        }
        \sdocarg{double ji}{
            Total angular momentum quantum number of the initial state
            $\bm{\Psi}_\mathrm{i}$, $j$. Here, the following convention is applied:
            For $j\in[k,k+1)$, $k=0,1,2,\dots$, the library internally uses the
            (half-integer) value $(2k+1)\slash{2}$. For instance, this means
            that for $j=0.49999$ and $j=0.50001$ the software internally uses
            $j=1\slash{2}$ in both cases.
        }
        \sdocarg{int32\tus{t} nf}{
            Principal quantum number of the final state $\bm{\Psi}_\mathrm{f}$.
        }
        \sdocarg{int32\tus{t} lf}{
            Orbital angular momentum quantum number $l=1,2,\dots,n-1$ of the
            final state $\bm{\Psi}_\mathrm{f}$.
        }
        \sdocarg{double sf}{
            Electron spin quantum number of the final state
            $\bm{\Psi}_\mathrm{f}$. Not an argument here, since $s=1\slash{2}$
            is always assumed.
        }
        \sdocarg{double jf}{
            Total angular momentum quantum number of the final state
            $\bm{\Psi}_\mathrm{f}$, $j$. Here, the following convention is
            applied: For $j\in[k,k+1)$, $k=0,1,2,\dots$, the library internally
            uses the (half-integer) value $(2k+1)\slash{2}$. For instance, this
            means that for $j=0.49999$ and $j=0.50001$ the software internally
            uses $j=1\slash{2}$ in both cases.
        }
        \sdocret{double fitof}{
            $f_\mathrm{i\to{f}}$ is the oscillator strength of the selected
            electronic transition, summed over $m'_j$ of the final state. Note
            that because of this sum, it is the same for all $m_j$ of this
            initial state.
        }
    \end{sdocfun}%
    %

    %%% BIBLIOGRAPHY
    \newpage
    \renewcommand\bibname{References}%
    \printbibliography

\end{document}
